\section{Introduction}

\label{sec:introduction}

The throughput of a single consensus instance is fundamentally limited by the per-round
message complexity and the finality latency. Even with pipelining (Photon) and burst
finalization (Nova), a single consensus instance processes transactions sequentially
within each conflict set. For blockchains with diverse transaction types---token
transfers, smart contract calls, cross-chain messages, governance votes---the aggregate
throughput is limited by the busiest conflict set.

Prism addresses this by partitioning the consensus problem into independent spectrum
bands, each handling a subset of conflict sets. This is analogous to frequency-division
multiplexing in telecommunications: each band occupies a different ``frequency''
(conflict set partition) and processes transactions independently.

The challenge is ensuring global consistency: while each band finalizes independently,
the combined output must form a valid total order compatible with cross-band dependencies.
Prism's spectral combiner solves this through a deterministic merge algorithm that
respects the causal order of cross-band references.

\paragraph{Contributions.}
\begin{enumerate}[nosep]
\item The Prism multi-spectrum consensus protocol with hash-based band assignment (\S\ref{sec:protocol}).
\item Conflict locality and load balance proofs (\S\ref{sec:partition}).
\item The spectral combiner algorithm with order consistency proof (\S\ref{sec:combiner}).
\item Direct voting fast path for uncontested decisions (\S\ref{sec:voting}).
\item Linear throughput scaling demonstration: 50,000 TPS with 16 bands (\S\ref{sec:evaluation}).
\end{enumerate}

%% -----------------------------------------------------------------------
%%  2. SPECTRAL DECOMPOSITION
%% -----------------------------------------------------------------------
